\section{Typical rank and spectral properties}\label{sec:typical}

We derive three properties of binary matrices that will be used below.
The rank properties hold simultaneously on all rational subspaces of
the prescribed height and for all choices of columns. Their uniformity
comes from Theorem~\ref{thm:low-height-cover}, together with a rank test
over \(\F_2\). The spectral property has the same exceptional-probability
scale, although its proof uses concentration rather than a union bound.

For a real symmetric matrix \(B\) and a subspace \(U\subseteq\R^m\),
write \(B|_U\) for the bilinear form
\((u,v)\mapsto u^{\mathsf T}Bv\) restricted to \(U\times U\).
Its nullity is \(\dim U-\rank(B|_U)\). We write \(\pi_U\) for
orthogonal projection onto \(U\), and put
\(E_X=\Span\{e_i:i\in X\}\) for \(X\subseteq[m]\).

\subsection{Uniform rank tests}

We first record the finite-field estimate behind both rank assertions.
Let \(\mathcal B\) be a uniform symmetric bilinear form on \(\F_2^m\),
including independent diagonal entries, and let \(S,T\) be fixed
subspaces of dimensions \(b,a\), respectively, where \(b\le a\).
For \(1\le t\le b\),
\begin{equation}\label{eq:typical-f2-rank}
 \Pr\bigl(\rank(\mathcal B|_{S\times T})\le b-t\bigr)
 \le C_0\,2^{-t(t+1)/2},
\end{equation}
with an absolute constant \(C_0\). Indeed, such a rank deficiency
provides a \(t\)-dimensional subspace \(K\subseteq S\) with
\(\mathcal B(K,T)=0\). The number of possible \(K\) is
\[
 \binom bt_{\!2}
 =2^{t(b-t)}\prod_{i=0}^{t-1}
       \frac{1-2^{i-b}}{1-2^{i-t}}
 \le C_0\,2^{t(b-t)},
 \qquad
 C_0=\prod_{j=1}^{\infty}(1-2^{-j})^{-1}<\infty.
\]
For a fixed \(K\), set \(e=\dim(K\cap T)\). Choose a basis of the
intersection, extend it separately to bases of \(K\) and \(T\), and
then extend their union to a basis of \(\F_2^m\). A change of basis
preserves the uniform distribution on symmetric forms. In this basis,
the equations \(\mathcal B(K,T)=0\) specify
\[
 ta-\binom e2
\]
independent entries of the symmetric form. In fact there are \(ta\)
entries before using symmetry, and the only repetitions are the two
orders of each off-diagonal pair inside \(K\cap T\). Diagonal entries
in that intersection remain independent conditions, also in
characteristic two. Since \(e\le t\) and \(a\ge b\), the number of
conditions is at least \(tb-\binom t2\). A union bound over \(K\)
gives Equation~\eqref{eq:typical-f2-rank}. The event is empty if
\(t>b\).

\begin{proposition}[Uniform low-height rank tests]\label{prop:typical-ranks}
For every pair of fixed constants \(D,A>0\), there is a function
\(f_{D,A}(N)\to\infty\) such that the following holds for every
\(m\in\mathcal I_N\) and all sufficiently large \(N\), with probability
at least \(1-2^{-Nf_{D,A}(N)}\) for \(B=\mathbf B_m\).
Simultaneously for every rational subspace \(W\subseteq\R^m\) with
\(H(W)\le\exp(N\ell^D)\), one has
\begin{enumerate}
 \item
 \(\dim W^\perp-\rank(B|_{W^\perp})\le N\ell^{-A}\);
 \item for every coordinate set \(X\subseteq[m]\),
 \[
  \rank\bigl[\pi_{W^\perp}Be_i\bigr]_{i\in X}
  \ge \min\{|X|,m-\dim W\}-N\ell^{-A}.
 \]
\end{enumerate}
The ranks and the orthogonal projections in these conclusions are
over \(\R\); no nonsingularity assumption on \(B\) is needed.
\end{proposition}

\begin{proof}
For any rational subspace \(U\subseteq\R^m\), the lattice
\(U\cap\Z^m\) is primitive in \(\Z^m\). Consequently an integer
basis matrix \(C_U\) of this lattice has full column rank modulo two:
a basis of a primitive sublattice extends to a basis of \(\Z^m\).
If \(U,V\) are rational, the matrix of the tested form is the integer
matrix \(C_U^{\mathsf T}BC_V\). Its real rank is at least the rank of
its reduction modulo two, because a nonzero minor modulo two is an
odd, hence nonzero, integer minor. Moreover \(B\bmod 2\) is uniform
among symmetric matrices over \(\F_2\). Thus
Equation~\eqref{eq:typical-f2-rank} applies to the reductions of the
column spaces of \(C_U,C_V\), which have the same respective
dimensions as \(U,V\). Their intersection after reduction can be
larger than the reduction of \(U\cap V\); the intersection correction
in the preceding count accommodates this without any further
assumption.

The tests we need are \(U=V\), and \(V=E_X\). In the second case,
\[
 \rank\bigl[\pi_UBe_i\bigr]_{i\in X}
   =\rank(C_U^{\mathsf T}BC_{E_X}),
\]
where \(C_{E_X}\) consists of the coordinate vectors with indices in
\(X\). To see this equality, \(C_U^{\mathsf T}\) annihilates
\(U^\perp\) and is an isomorphism from \(U\) to \(\R^{\dim U}\).
It therefore detects exactly the rank after projection to \(U\).

We also need stability when the testing subspace changes. For fixed
\(B\) and \(X\), define
\[
 \delta_0(U)=\dim U-\rank(B|_U),\qquad
 \delta_X(U)=\min\{|X|,\dim U\}
      -\rank(C_U^{\mathsf T}BC_{E_X}),
\]
where the second rank can equivalently be defined by the bilinear
form on \(U\times E_X\), without choosing an integer basis. If
\(Z\subseteq U\), restricting a matrix to \(Z\times Z\) deletes
\(\dim U-\dim Z\) rows and the same number of columns. Its rank
decreases by an amount between zero and twice that number. Hence
\[
 |\delta_0(U)-\delta_0(Z)|\le\dim U-\dim Z.
\]
For the form on \(U\times E_X\), the rank and the quantity
\(\min\{|X|,\dim U\}\) each decrease by an amount between zero and
\(\dim U-\dim Z\). The same inequality follows for \(\delta_X\).
Taking \(Z=U\cap V\) in these two inequalities gives
\begin{equation}\label{eq:typical-rank-stability}
 |\delta_0(U)-\delta_0(V)|\le d(U,V),\qquad
 |\delta_X(U)-\delta_X(V)|\le d(U,V).
\end{equation}
Here \(d(U,V)=\dim U+\dim V-2\dim(U\cap V)\), and the dimension
formula also gives \(d(U^\perp,V^\perp)=d(U,V)\).

Fix \(D,A\), put \(s=N\ell^{-A}\), and apply
Theorem~\ref{thm:low-height-cover} with precision exponent
\(P=2A+4\). For each \(m\in\mathcal I_N\), it gives a list
\(\mathcal C_m\) of rational subspaces with
\[
 |\mathcal C_m|\le 2^{N^2\ell^{-P}},\qquad
 \min_{V\in\mathcal C_m}d(W,V)\le N\ell^{-P}
\]
for every \(W\) in the height range of the proposition. Set
\(t=\lceil s/2\rceil\). Apply Equation~\eqref{eq:typical-f2-rank}
to \(U=V^\perp\) for every \(V\in\mathcal C_m\), both to its
restricted form and to its form against every \(E_X\). The comparison
of real and mod-two ranks above shows that the probability that any
of these real deficiencies is at least \(t\) is at most
\[
 C_0(1+2^m)\,2^{N^2\ell^{-P}-t(t+1)/2}
 \le 2^{-\Omega(N^2\ell^{-2A})}.
\]
In the last inequality, \(N^2\ell^{-P}=o(N^2\ell^{-2A})\) and
\(m=o(N^2\ell^{-2A})\), uniformly in the stated range of \(m\).
Tests whose smaller dimension is less than \(t\) need not be included.

Outside this event, choose a covering \(V\) for a given \(W\).
Equation~\eqref{eq:typical-rank-stability}, applied to their
orthogonal complements, bounds each deficiency for \(W\) by
\[
 t-1+N\ell^{-P}<s
\]
for all sufficiently large \(N\). This holds simultaneously for
all such \(W\) and all \(X\). Finally,
\(N\ell^{-2A}\to\infty\), so the displayed exceptional probability
has the form required in the statement.
\end{proof}

\begin{corollary}[Inverse restriction rank]\label{cor:inverse-rank}
For every pair of fixed constants \(D,A>0\), typically for nonsingular
binary matrices \(B\) of size \(m\in\mathcal I_N\), simultaneously
for every rational \(W\subseteq\R^m\) with
\(H(W)\le\exp(N\ell^D)\),
\[
 \rank(B^{-1}|_W)\ge\dim W-N\ell^{-A}.
\]
Here \(B^{-1}|_W\) is the restriction of the bilinear form with matrix
\(B^{-1}\), rather than the restriction of its associated linear map.
The exceptional probability is uniform in \(m\), as in
Proposition~\ref{prop:typical-ranks}.
\end{corollary}

\begin{proof}
The two radicals are
\[
 \operatorname{rad}(B|_{W^\perp})
   =W^\perp\cap B^{-1}W,
 \qquad
 \operatorname{rad}(B^{-1}|_W)
   =W\cap B(W^\perp).
\]
Since \(B\) is nonsingular, multiplication by \(B\) is an
isomorphism from the first space onto the second, with inverse
multiplication by \(B^{-1}\). Their dimensions are equal. The first
assertion of Proposition~\ref{prop:typical-ranks} therefore gives the
claimed rank for the inverse form, on the same typical set.
\end{proof}

\subsection{Corank modulo all primes}

For nonsingular \(B\) and every prime \(p\), the presentation of \(G_B\) gives
\[
 k_B(p,1)=\dim_{\F_p}(G_B/pG_B)
          =\corank_{\F_p}(B\bmod p).
\]
This dimension is the minimum number of generators of the \(p\)-primary
part, and \(k_B(p,e)\le k_B(p,1)\) for every \(e\). A uniform bound on
these coranks will therefore control the generators needed for quotients
of \(G_B\) and cap its layer heights.

\begin{proposition}[Uniform prime corank]\label{prop:prime-corank}
Typically for nonsingular binary matrices \(B\) of size
\(m\in\mathcal I_N\), one has, simultaneously for every prime \(p\),
\[
 \corank_{\F_p}(B\bmod p)\le N^{3/5}.
\]
More precisely, uniformly for \(m\in\mathcal I_N\), the probability
of nonsingularity together with failure of this assertion is at most
\(2^{-\Omega(N^{6/5})}\).
\end{proposition}

\begin{proof}
If \(B\) is nonsingular over \(\R\), its determinant is a nonzero
integer. Every prime at which its reduction has positive corank
divides that determinant. Hadamard's inequality gives
\[
 1\le |\det B|\le m^{m/2}\le N^{N/2},
\]
so it suffices to consider primes \(p\le N^{N/2}\). For each such
prime, Lemma~\ref{lem:corank-tail}, with field \(\F_p\), bounds the
probability of corank greater than \(N^{3/5}\) by
\(2^{-cN^{6/5}}\) for an absolute \(c>0\), uniformly in \(p\) and
\(m\). Even bounding the number of these primes by \(N^{N/2}\),
the union bound is
\[
 2^{(N/2)\log_2N-cN^{6/5}}
   \le 2^{-cN^{6/5}/2}
\]
for all sufficiently large \(N\). This also proves the asserted
uniformity.
\end{proof}

\subsection{Few eigenvalues near zero}

\begin{proposition}[Typical small-spectrum bound]\label{prop:typical-spectrum}
For every fixed \(A>0\), there is a fixed \(C_A>0\) and a function
\(f_A(N)\to\infty\) such that, for every \(m\in\mathcal I_N\) and
all sufficiently large \(N\), with probability at least
\(1-2^{-Nf_A(N)}\), the number of eigenvalues of \(B=\mathbf B_m\)
in
\[
 [-\sqrt N\,\ell^{-C_A},\ \sqrt N\,\ell^{-C_A}]
\]
is at most \(N\ell^{-A}\). Eigenvalues are counted with multiplicity.
One may take \(C_A=A+6\). This assertion does not require \(B\) to
be nonsingular.
\end{proposition}

\begin{proof}
\textbf{Mean bound.}
Let \(J_m\) be the all-ones matrix and write
\[
 2B=J_m+S,\qquad X=S/\sqrt m.
\]
The entries of \(S\) on and above the diagonal are independent uniform
signs. We first bound the number of eigenvalues of \(X\) in a short
interval about zero.

The limiting moments in the following calculation are those of
Wigner's semicircle law~\cite{Wigner}. We include the closed-walk
enumeration to obtain the quantitative error bound needed here. Let
\[
 d\mu_{\mathrm{sc}}(x)=\frac{1}{2\pi}\sqrt{4-x^2}\,
             \ind_{[-2,2]}(x)\,dx.
\]
For \(i\ge1\) with \(i+1\le m\),
\begin{equation}\label{eq:typical-walk-moments}
 \left|\frac1m\E\Tr X^{2i}
       -\int x^{2i}\,d\mu_{\mathrm{sc}}(x)\right|
 \le \frac{(2i+1)^{2i+2}}{m},
 \qquad
 \E\Tr X^{2i-1}=0.
\end{equation}
To prove it, expand the normalized trace as
\[
 \frac1m\E\Tr X^k
 =m^{-1-k/2}\sum_{v_0,\ldots,v_{k-1}\in[m]}
       \E\prod_{r=0}^{k-1}S_{v_rv_{r+1}},
 \qquad v_k=v_0.
\]
Each summand is a closed walk, with edges regarded as unordered and
loops allowed. Its expectation is one if every distinct edge is used
an even number of times, and is zero otherwise. This proves the
odd-moment assertion. For \(k=2i\), a contributing walk has at most
\(i\) distinct edges. Its graph is connected, and hence has at most
\(i+1\) vertices. Equality forces exactly \(i\) nonloop edges,
each used twice, forming a tree. When vertices are named in order of
first appearance, these walks are in bijection with Dyck words of
length \(2i\): the first traversal of a tree edge raises the depth,
and its return traversal lowers it. Their number is the Catalan number
\(\mathrm{Cat}_i=(i+1)^{-1}\binom{2i}{i}\), and each has
\((m)_{i+1}=m(m-1)\cdots(m-i)\) injective labelings.

The contribution of these tree walks is therefore
\(\mathrm{Cat}_i(m)_{i+1}/m^{i+1}\). Its difference from
\(\mathrm{Cat}_i\) is at most
\(\mathrm{Cat}_i i(i+1)/(2m)\), using
\(1-\prod_{h=0}^i(1-h/m)\le\sum_{h=0}^i h/m\).
Every remaining contributing walk has at most \(i\) vertices.
There are at most \((2i+1)^{2i}\) canonical vertex sequences, a
crude bound obtained by allowing each position any of \(2i+1\)
names. Each such sequence has at most \(m^i\) labelings and so
contributes at most \(1/m\) after normalization. Finally, elementary
beta integration gives
\[
 \int x^{2i}\,d\mu_{\mathrm{sc}}(x)
  =\frac{(2i)!}{i!(i+1)!}=\mathrm{Cat}_i,
\]
and its odd moments vanish. Combining these estimates, and using
\(\mathrm{Cat}_i\le4^i\), gives
Equation~\eqref{eq:typical-walk-moments}.

Fix the exponent \(A\) in the proposition and put
\[
 w=\ell^{-A-5},\qquad j=\lceil\ell^{2A+4}\rceil,
 \qquad f(x)=(1-|x|/w)_+,
 \qquad p_j(x)=2(1-x^2/16)^{2j}.
\]
For all large \(N\), \(f(x)\le p_j(x)\) for every real \(x\).
Outside \([-w,w]\) this follows from nonnegativity of the even
power. Inside that interval, Bernoulli's inequality gives
\[
 p_j(x)\ge 2(1-w^2/16)^{2j}
       \ge 2(1-jw^2/8)\ge1\ge f(x),
\]
because \(jw^2=O(\ell^{-6})\). All moments in \(p_j\) have degree
at most \(4j\), and \(2j+1\le m\) uniformly in
\(m\in\mathcal I_N\) for large \(N\). Expanding this polynomial
and applying Equation~\eqref{eq:typical-walk-moments} term by term,
the sum of the absolute values of its coefficients is
\(2(17/16)^{2j}\). Consequently
\[
 \left|\frac1m\E\Tr p_j(X)
        -\int p_j\,d\mu_{\mathrm{sc}}\right|
 \le \frac{2(17/16)^{2j}(4j+1)^{4j+2}}{m}
 =m^{-1+o(1)}.
\]
This estimate verifies the required growing degrees: indeed
\[
 j\log(j+1)=O_A(\ell^{2A+4}\log\ell)=o(\log m)
\]
uniformly in the stated range. In particular the last error is
\(o(\ell^{-A})\). On the support of \(\mu_{\mathrm{sc}}\),
\((1-x^2/16)^{2j}\le\exp(-jx^2/8)\) and the density is at most
\(1/\pi\). Thus
\[
 \int p_j\,d\mu_{\mathrm{sc}}
 \le \frac2\pi\int_{\R}e^{-jx^2/8}\,dx
 =O(j^{-1/2})=O(\ell^{-A-2}).
\]
The pointwise majorization now proves
\begin{equation}\label{eq:typical-tent-mean}
 \E\Tr f(X)=o(m\ell^{-A}),
\end{equation}
again uniformly in \(m\).

\medskip\noindent\textbf{Concentration and translation to \(B\).}
We next obtain concentration strong enough for typicality. Pointwise,
\[
 f(x)=1-g_1(x)+g_2(x),\qquad
 g_1(x)=|x|/w,\qquad g_2(x)=(|x|/w-1)_+.
\]
Both \(g_1,g_2\) are convex and \(1/w\)-Lipschitz. Their traces are
convex functions of a real symmetric matrix. For completeness, if
\(M_+\) denotes its spectral positive part, then
\[
 \Tr M_+=\max_{0\preceq P\preceq I}\Tr(PM).
\]
The identities
\[
 \Tr g_1(M)=w^{-1}\bigl(\Tr M_++\Tr(-M)_+\bigr),\qquad
 \Tr g_2(M)=\Tr(M/w-I)_++\Tr(-M/w-I)_+
\]
express each trace as a sum of suprema of affine functions of \(M\),
which proves convexity.

Let \(y=(y_{ab})_{a\le b}\) be an arbitrary real vector of independent
upper-triangular coordinates, and form \(X(y)=S(y)/\sqrt m\) by
symmetric reflection. For real symmetric matrices, the
Hoffman--Wielandt inequality~\cite[Theorem~1 and Remark~2]{HoffmanWielandt}, with their
eigenvalues ordered, states that
\[
 \sum_{i=1}^m|\lambda_i(M)-\lambda_i(M')|^2
       \le\|M-M'\|_F^2.
\]
It follows by Cauchy--Schwarz that for \(a=1,2\),
\[
 |\Tr g_a(X(y))-\Tr g_a(X(z))|
 \le\frac{\sqrt m}{w}\|X(y)-X(z)\|_F
 \le\frac{\sqrt2}{w}\|y-z\|_2.
\]
The last inequality uses
\(\|X(y)-X(z)\|_F^2\le(2/m)\|y-z\|_2^2\).
Thus each trace is convex and \(L\)-Lipschitz as a function of the
independent coordinates, with \(L=\sqrt2/w\).

Talagrand's convex Lipschitz concentration theorem on the product
cube~\cite[Theorem~3]{Talagrand} implies that, for universal \(c,C>0\),
if \(F:\R^d\to\R\) is convex and globally \(L\)-Lipschitz for the
Euclidean norm, then its value at independent uniform signs satisfies
\[
 \Pr(|F-\E F|>u)\le C\exp(-cu^2/L^2)\qquad(u\ge0).
\]
The usual statement about a median implies this form: integrating
its tail bounds the difference between the median and the expectation
by \(O(L)\), which can be absorbed by adjusting \(c,C\).
Apply the theorem to the two traces. By
Equation~\eqref{eq:typical-tent-mean}, their difference
\(\Tr f(X)=m-\Tr g_1(X)+\Tr g_2(X)\) has expectation at most
\(m\ell^{-A}/8\) for large \(N\). If it exceeds
\(m\ell^{-A}/4\), at least one of the two convex traces deviates
from its expectation by more than \(m\ell^{-A}/16\). Therefore
\begin{equation}\label{eq:typical-tent-tail}
 \Pr\bigl(\Tr f(X)>m\ell^{-A}/4\bigr)
 \le 2C\exp\bigl(-c' m^2\ell^{-2A}w^2\bigr)
 \le 2^{-c_A N^2\ell^{-4A-10}}
\end{equation}
for some \(c',c_A>0\) and all sufficiently large \(N\), uniformly
in \(m\). Since \(N\ell^{-4A-10}\to\infty\), this is a
\(2^{-Nf_A(N)}\) exceptional probability with \(f_A(N)\to\infty\).

On the complementary event, \(f\ge1/2\) on \([-w/2,w/2]\) shows
that \(X\) has at most \(m\ell^{-A}/2\) eigenvalues in that
interval. Finally \(2B/\sqrt m=X+J_m/\sqrt m\) is a rank-one
perturbation of \(X\). Eigenvalue interlacing changes the cumulative
eigenvalue count at either endpoint of an interval by at most one,
so it can increase the count in any interval by at most two. Take
\(C_A=A+6\). Uniformly for \(m\in\mathcal I_N\),
\[
 2\sqrt{N/m}\,\ell^{-C_A}\le w/2
\]
for all sufficiently large \(N\). The eigenvalues of \(B\) in the
interval of the proposition therefore number at most
\[
 \frac12m\ell^{-A}+2\le N\ell^{-A}.
\]
Together with Equation~\eqref{eq:typical-tent-tail}, this proves the
proposition.
\end{proof}
